% GENERATED FILE. Edit canonical lessons, then run npm run generate:books.
% canonical-source-hash: fnv1a64:2cfebc27b514ac71
% canonical-lessons: FA-C530-dastan-e-kutah, FA-C530-shakhsiyat-e-asli, FA-C530-tarsnak, FA-C530-ghamangiz, FA-C530-khandedar

\chapter{Kinds of Stories}
\label{ch:fa-kinds-of-stories}
\hlchaptermodality{530}

\begin{chapteropening}
\textbf{By the end of this chapter:} \emph{I can say what kind of story something is.}

Complete the last lesson of chapter 530: I can say what kind of story something is.
\end{chapteropening}

\section[\texorpdfstring{\fa{داستان} \fa{کوتاه} (dâstân-e kutâh)}{داستان کوتاه (dâstân-e kutâh)}]{\fa{داستان} \fa{کوتاه} (dâstân-e kutâh) — a short story}

\label{lesson:FA-C530-dastan-e-kutah}



\begin{quote}
Before the new one: say the Persian for to study, then the Persian for a bookworm.
\end{quote}

\subsection*{You'll want to know: \fa{داستان} \fa{کوتاه}}
\textbf{\fa{داستان} \fa{کوتاه}} (\emph{dâstân-e kutâh}) — \textquotedblleft{}a short story\textquotedblright{}.

\begin{quote}
\fa{این} \fa{نویسنده} \fa{بیشتر} \fa{داستان} \fa{کوتاه} \fa{مینویسد}.
\end{quote}

\begin{quote}
\emph{in nevisande bishtar dâstân-e kutâh minevisad.}
\end{quote}

\begin{quote}
This writer mostly writes short stories.
\end{quote}

\emph{dâstân} 'story' + \emph{kutâh} 'short'. Plural \emph{dâstânhâ-ye kutâh}.

\begin{cousinweb}
\emph{dâstân-e kutâh} is built on two native Persian words, \emph{dâstân}, 'story', and \emph{kutâh}, 'short', which EtymDB sets beside \emph{kuchek}, 'small' (EtymDB).
\end{cousinweb}

\begin{grammarlens}[title={What you've built}]
One more word for telling, explaining and describing at B1.
\end{grammarlens}

\subsection*{Your turn}
\begin{itemize}
\raggedright
  \item \emph{Say it:} \emph{dâstân-e kutâh}
  \item \emph{Say it:} \emph{dâstân-e kutâh}, once more
  \item \emph{Say it:} a sentence of your own with \emph{dâstân-e kutâh}
  \item \emph{Recall:} say the Persian for to study, then the Persian for a bookworm, then say \emph{dâstân-e kutâh} again
\end{itemize}

\begin{tcolorbox}[breakable,colback=teal!4,colframe=teal!35!black,title={Before you move on}]
What does \fa{داستان} \fa{کوتاه} mean? (\textquotedblleft{}A short story\textquotedblright{}.) Say it once more.
\end{tcolorbox}
\section[\texorpdfstring{\fa{شخصیت} \fa{اصلی} (shakhsiyat-e asli)}{شخصیت اصلی (shakhsiyat-e asli)}]{\fa{شخصیت} \fa{اصلی} (shakhsiyat-e asli) — the main character}

\label{lesson:FA-C530-shakhsiyat-e-asli}



\begin{quote}
Before the new one: say the Persian for a bookworm, then the Persian for a short story.
\end{quote}

\subsection*{You'll want to know: \fa{شخصیت} \fa{اصلی}}
\textbf{\fa{شخصیت} \fa{اصلی}} (\emph{shakhsiyat-e asli}) — \textquotedblleft{}the main character\textquotedblright{}.

\begin{quote}
\fa{شخصیت} \fa{اصلی} \fa{این} \fa{رمان} \fa{یک} \fa{دختر} \fa{جوان} \fa{است}.
\end{quote}

\begin{quote}
\emph{shakhsiyat-e asli-ye in român yek dokhtar-e javân ast.}
\end{quote}

\begin{quote}
The main character of this novel is a young girl.
\end{quote}

\emph{shakhsiyat} 'character, personality' + \emph{asli} 'main'.

\begin{cousinweb}
\emph{shakhsiyat-e asli} is built in Persian from two Arabic loans: \emph{shakhsiyat}, from Arabic \emph{šaḫṣiyya} (EtymDB), and \emph{asli}, from Arabic \emph{aṣlī} (EtymDB).
\end{cousinweb}

\begin{grammarlens}[title={What you've built}]
One more word for telling, explaining and describing at B1.
\end{grammarlens}

\subsection*{Your turn}
\begin{itemize}
\raggedright
  \item \emph{Say it:} \emph{shakhsiyat-e asli}
  \item \emph{Say it:} \emph{shakhsiyat-e asli}, once more
  \item \emph{Say it:} a sentence of your own with \emph{shakhsiyat-e asli}
  \item \emph{Recall:} say the Persian for a bookworm, then the Persian for a short story, then say \emph{shakhsiyat-e asli} again
\end{itemize}

\begin{tcolorbox}[breakable,colback=teal!4,colframe=teal!35!black,title={Before you move on}]
What does \fa{شخصیت} \fa{اصلی} mean? (\textquotedblleft{}The main character\textquotedblright{}.) Say it once more.
\end{tcolorbox}
\section[\texorpdfstring{\fa{ترسناک} (tarsnâk)}{ترسناک (tarsnâk)}]{\fa{ترسناک} (tarsnâk) — scary, frightening}

\label{lesson:FA-C530-tarsnak}



\begin{quote}
Before the new one: say the Persian for a short story, then the Persian for the main character.
\end{quote}

\subsection*{You'll want to know: \fa{ترسناک}}
\textbf{\fa{ترسناک}} (\emph{tarsnâk}) — \textquotedblleft{}scary, frightening\textquotedblright{}.

\begin{quote}
\fa{دیشب} \fa{یک} \fa{فیلم} \fa{ترسناک} \fa{دیدیم}.
\end{quote}

\begin{quote}
\emph{dishab yek film-e tarsnâk didim.}
\end{quote}

\begin{quote}
Last night we watched a scary film.
\end{quote}

\emph{tars} 'fear' + \emph{-nâk} 'full of'.

\begin{cousinweb}
\emph{tarsnâk} is built on native Persian \emph{tars}, 'fear', from inherited \emph{tarsidan}, 'to fear', with the ending \emph{-nâk}, 'full of'.
\end{cousinweb}

\begin{grammarlens}[title={What you've built}]
One more word for telling, explaining and describing at B1.
\end{grammarlens}

\subsection*{Your turn}
\begin{itemize}
\raggedright
  \item \emph{Say it:} \emph{tarsnâk}
  \item \emph{Say it:} \emph{tarsnâk}, once more
  \item \emph{Say it:} a sentence of your own with \emph{tarsnâk}
  \item \emph{Recall:} say the Persian for a short story, then the Persian for the main character, then say \emph{tarsnâk} again
\end{itemize}

\begin{tcolorbox}[breakable,colback=teal!4,colframe=teal!35!black,title={Before you move on}]
What does \fa{ترسناک} mean? (\textquotedblleft{}Scary, frightening\textquotedblright{}.) Say it once more.
\end{tcolorbox}
\section[\texorpdfstring{\fa{غم} \fa{انگیز} (ghamangiz)}{غم انگیز (ghamangiz)}]{\fa{غم} \fa{انگیز} (ghamangiz) — tragic, sad}

\label{lesson:FA-C530-ghamangiz}



\begin{quote}
Before the new one: say the Persian for the main character, then the Persian for scary.
\end{quote}

\subsection*{You'll want to know: \fa{غم} \fa{انگیز}}
\textbf{\fa{غم} \fa{انگیز}} (\emph{ghamangiz}) — \textquotedblleft{}tragic, sad\textquotedblright{}.

\begin{quote}
\fa{پایان} \fa{داستان} \fa{غم} \fa{انگیز} \fa{بود}.
\end{quote}

\begin{quote}
\emph{pâyân-e dâstân ghamangiz bud.}
\end{quote}

\begin{quote}
The end of the story was sad.
\end{quote}

\emph{gham} 'sorrow' + \emph{angiz} 'stirring' (stem of \emph{angikhtan}): sorrow-stirring.

\begin{cousinweb}
\emph{ghamangiz} is built in Persian on \emph{gham}, 'sorrow', from Arabic \emph{ġamm} (EtymDB), and \emph{angiz}, the stem of native \emph{angikhtan}, 'to stir up'.
\end{cousinweb}

\begin{grammarlens}[title={What you've built}]
One more word for telling, explaining and describing at B1.
\end{grammarlens}

\subsection*{Your turn}
\begin{itemize}
\raggedright
  \item \emph{Say it:} \emph{ghamangiz}
  \item \emph{Say it:} \emph{ghamangiz}, once more
  \item \emph{Say it:} a sentence of your own with \emph{ghamangiz}
  \item \emph{Recall:} say the Persian for the main character, then the Persian for scary, then say \emph{ghamangiz} again
\end{itemize}

\begin{tcolorbox}[breakable,colback=teal!4,colframe=teal!35!black,title={Before you move on}]
What does \fa{غم} \fa{انگیز} mean? (\textquotedblleft{}Tragic, sad\textquotedblright{}.) Say it once more.
\end{tcolorbox}
\section[\texorpdfstring{\fa{خنده} \fa{دار} (khandedâr)}{خنده دار (khandedâr)}]{\fa{خنده} \fa{دار} (khandedâr) — funny}

\label{lesson:FA-C530-khandedar}



\begin{quote}
Before the new one: say the Persian for scary, then the Persian for tragic.
\end{quote}

\subsection*{You'll want to know: \fa{خنده} \fa{دار}}
\textbf{\fa{خنده} \fa{دار}} (\emph{khandedâr}) — \textquotedblleft{}funny\textquotedblright{}.

\begin{quote}
\fa{این} \fa{فیلم} \fa{خیلی} \fa{خنده} \fa{دار} \fa{بود}.
\end{quote}

\begin{quote}
\emph{in film kheyli khandedâr bud.}
\end{quote}

\begin{quote}
This film was very funny.
\end{quote}

\emph{khande} 'laughter' + \emph{dâr} 'having'. \emph{bâmaze} is 'cute, witty'.

\begin{cousinweb}
\emph{khandedâr} is built on native Persian \emph{khande}, 'laughter', from \emph{khandidan}, 'to laugh', and \emph{dâr}, the stem of \emph{dâshtan}, 'to have'.
\end{cousinweb}

\begin{grammarlens}[title={What you've built}]
One more word for telling, explaining and describing at B1.
\end{grammarlens}

\subsection*{Your turn}
\begin{itemize}
\raggedright
  \item \emph{Say it:} \emph{khandedâr}
  \item \emph{Say it:} \emph{khandedâr}, once more
  \item \emph{Say it:} a sentence of your own with \emph{khandedâr}
  \item \emph{Recall:} say the Persian for scary, then the Persian for tragic, then say \emph{khandedâr} again
\end{itemize}

\begin{tcolorbox}[breakable,colback=teal!4,colframe=teal!35!black,title={Before you move on}]
What does \fa{خنده} \fa{دار} mean? (\textquotedblleft{}Funny\textquotedblright{}.) Say it once more.
\end{tcolorbox}
